% !PRINTED 559
\weberpage{559}
\begin{center}
\gesperrt{Achtzehnter Abschnitt.}

\gesperrt{Discriminanten.}

\rule{0.3\textwidth}{0.4pt}

\S.\,155.

\gesperrt{Minimum einer quadratischen Form.}
\end{center}

Auf einem ganz neuen Weg hat \gesperrt{Minkowski} die Endlichkeit der Anzahl der Classen bewiesen\footnote{Ueber die positiven quadratischen Formen etc. \gesperrt{Crelle}, Bd.\ 107, 1891. Eingehender noch in dem Buche von \gesperrt{Minkowski}, \glqq{}Geometrie der Zahlen\grqq{}, Leipzig 1896.}, der von eigenthümlichen, zahlreicher Anwendungen fähigen Betrachtungen ausgeht und dabei zu dem lange vergeblich gesuchten Beweis eines Satzes über die Körperdiscriminante führt. Wir wollen hier in der Kürze die Hauptmomente dieser Untersuchungen mittheilen.

\gesperrt{Minkowski} geht aus von der Betrachtung einer quadratischen Form von $n$ Variablen

\weq{1}{f(x_1, x_2, \ldots, x_n) = \Sigma^{i,k} a_{i,k}\, x_i x_k}

mit reellen Co\"efficienten $a_{i,k}$, von der vorausgesetzt wird, dass sie sich in $n$ und nicht weniger positive Quadrate linearer Functionen zerlegen lässt (einer positiven Form, Bd.\ I, \S.\,81 bis 83). Es möge eine dieser Darstellungen sein

\weq{2}{f(x_1, x_2, \ldots, x_n) = \xi_1^2 + \xi_2^2 + \cdots + \xi_n^2,}
worin
\weq{3}{\xi_i = \alpha_{1,i} x_1 + \alpha_{2,i} x_2 + \cdots + \alpha_{n,i} x_n}

ein System linearer Functionen der $x_i$ mit nicht verschwindender Determinante ist. Setzen wir

\weq{4}{A = \Sigma \pm \alpha_{1,1}\alpha_{2,2}\ldots\alpha_{n,n},}
%CONT>>
